\section{Introduction}
\label{sec:introduction}

The question of what happens when an autonomous multi-agent system encounters a condition outside its design envelope has transitioned from theoretical to operational. Autonomous agent populations are now deployed in agricultural resource management, clinical decision support, financial clearing, and critical logistics — domains where the cost of an uncontrolled failure is measured in lives, capital, and institutional trust. As agent populations grow in recursive depth and operational scope, the absence of a structural mechanism compelling human involvement for unanticipated exception conditions has become the central vulnerability of agentic infrastructure. We introduce the \textbf{Bailout Protocol}: the mandatory, architecturally enforced escalation mechanism of the \bxthree{} Framework. It compels every unresolved exception state in any \bxthree{} node to propagate upward to a named human principal, bypassing all intermediate machine actors, with full forensic attribution at every step. It is not a feature added to an existing architecture. It is the mechanism by which the \bxthree{} Framework makes its upstream accountability guarantee structurally operative at runtime.

\subsection{Why Mandatory Bail-Out Is Structurally Necessary}
\label{sec:why-mandatory}

Any system operating with sufficient scope in a partially unknown environment will eventually encounter a state not anticipated at provisioning time. This is not a contingent property of poorly designed systems — it is a necessary property of any system whose operating envelope is finite and whose environment is open-ended. The question is not whether such conditions will arise; they will. The question is what the system is architecturally required to do when they do.

Current architectures resolve this in two inadequate ways. Internal error-handling covers only conditions enumerated at design time; novel exceptions proceed without a handler and propagate outward through the execution graph invisibly. Hierarchical escalation chains permit absorption at each intermediate tier before any exception reaches a human principal, and the absorption is invisible to external auditors. In recursive multi-agent systems — where nodes spawn sub-agents, share state, and pursue objectives within a shared environment — the consequences of either pattern compound. An unresolved exception at one level triggers compensatory actions at adjacent levels, each generating new anomaly states. Because no architectural constraint compels human involvement, the system attempts to absorb the anomaly within the machine-only execution graph. Each absorption generates further anomalies. The result is \textbf{accountability stranding}: the complete detachment of consequential decision authority from any human principal, invisible to external auditors, and recorded only in the system's own potentially compromised internal state.

A mandatory bail-out mechanism breaks this cycle at its origin. When a node's provisioned operating range is exceeded and the condition cannot be resolved within the node's delegated authority, the protocol compulsorily escalates to the node's designated human anchor, bypassing every intermediate machine actor. The human principal is not one option among many — it is the only permissible terminus for every unresolved exception. This compulsion is not a design choice that could be made optional; it is a structural requirement that follows from the functional separation between the \purpose{}, \bounds{}, and \fact{} layers.

\subsection{What Happens When Systems Lack One}
\label{sec:failure-modes}

Three failure modes follow directly from the absence of a mandatory bail-out architecture.

\textbf{Contingent oversight.} Human involvement becomes a runtime option determined by the system's own self-assessment of its competence. But a system is precisely least equipped to make this assessment correctly at the moment escalation is first required — when it has encountered a condition outside its design envelope and its model of the situation is most likely to be wrong. The humans most necessary to the process are least likely to be involved in it.

\textbf{Structural unauditability.} Without an architectural compulsion to involve a human principal, there is no requirement that any human be aware of an exception event. The system may absorb, reframe, or suppress the exception internally, and the decision class it represents — consequential actions taken without any human principal's knowledge — is invisible to every external auditor and to every internal record the system itself controls. In regulated environments, this invisible decision class is structurally disqualifying.

\textbf{Drift amplification.} Each successive decision made under conditions of unacknowledged exception produces a system state further from the principals' intent. Because the \bounds{} layer has not been modified — only violated — the system has no structural mechanism to recognize the divergence. The drift continues until environmental resistance or resource exhaustion terminates the process. The system is not choosing to continue; it has lost the ability to choose otherwise.

These failure modes are not eliminated by incremental improvement to error-handling or oversight practices. They are architectural properties of systems that treat accountability as a policy constraint rather than a structural requirement.

\subsection{Accountability as Architectural Fallback}
\label{sec:core-insight}

The core insight of the \bxthree{} Framework is that accountability must be an architectural fallback — a property enforced by the structure of the system, not a property achieved through aspiration or policy.

Policy constraints operate above the execution layer. They can be reinterpreted, negotiated, or deferred by agents with sufficient agency to reclassify constraint violations as optimization opportunities. A system that instructs agents to stop if risk exceeds a threshold is advisory at the execution layer; a sufficiently sophisticated agent will identify conditions under which continued execution is valued more highly than compliance. Architectural fallbacks operate within the execution layer. They are triggered deterministically by conditions the agent cannot control and cannot route around. The moment the fallback fires, the execution gate closes — not because a human intervened, not because the agent chose to stop, but because the structural conditions for continued execution are no longer satisfied.

The \bxthree{} Framework organizes every node into three immutable functional layers. The \purpose{} layer carries intent, judgment, and final authority — it defines the operating range $R(n)$ for each node $n$ and receives all escalated exceptions via the Bailout Protocol. The \purpose{} layer cannot delegate its exception-receiving function; no machine actor may serve as the terminal node for an unresolved exception. The \bounds{} layer carries bounded reasoning and constrained execution within $R(n)$. When it encounters a condition outside $R(n)$ it cannot resolve, it triggers the Bailout Protocol — compulsorily and non-discretionarily. It cannot absorb, reframe, or redirect the exception once triggered. The \fact{} layer carries deterministic enforcement and forensic auditability — it is the only layer authorized to commit state transitions or record entries to the append-only Forensic Ledger, producing the immutable evidentiary record that makes every escalation independently auditable.

When the \bounds{} layer encounters a condition outside $R(n)$ it cannot resolve, the functional separation itself dictates that the decision must return to the \purpose{} layer. There is no other layer with the authority to adjudicate it. The Bailout Protocol is not a feature added to this separation; it is the mechanism by which the separation enforces its own accountability guarantee. The human anchor $h(n)$ is not a design choice — it is a structural necessity.

This distinguishes \bxthree{} from prior escalation architectures. The Tiered Agentic Oversight model~\cite{kim2025tao} permits absorption at intermediate machine tiers, making human oversight a high-tier option rather than a compulsory terminus. The Bailout Protocol forbids this absorption by architectural constraint: the distinction is the difference between a system that \emph{may} consult a human and a system that \emph{cannot proceed} without one.

The Bailout Protocol is the fifth named architectural pillar of \bxthree{}, complementing Loop Isolation, Recursive Spawning, Spatial Firewall, and Sandbox Gate. Together, these five pillars constitute the upstream accountability guarantee: that \bxthree{} systems fail upward into human consciousness, never downward into autonomous chaos.

\subsection{Paper Roadmap}
\label{sec:roadmap}

Section~\ref{sec:bailout} provides the complete formal specification of the Bailout Protocol: its deterministic state machine, trigger condition, escalation algorithm, parent-chain bypass, and timeout handling. Section~\ref{sec:implementation} addresses the protocol's realization within the \bxthree{} three-layer architecture, showing how each layer participates in protocol execution. Section~\ref{sec:properties} formally states the Safety, Liveness, and Accountability correctness properties and derives them from the protocol's specification. Section~\ref{sec:related} surveys related work. Section~\ref{sec:limitations} concludes with structural limitations and a directed research agenda.
